\section{Current Standing Under the 2024 Norms}

\subsection{The current procedural law}

Since Pentecost 2024 the discernment of alleged supernatural
phenomena is governed by the Dicastery for the Doctrine of the
Faith's \work{Norms for Proceeding in the Discernment of Alleged
Supernatural Phenomena}, given 17 May 2024, approved by Pope
Francis on 4 May 2024, effective 19 May 2024, replacing the norms
of 25 February 1978. Under them the diocesan bishop examines
cases in his territory and formulates a judgment ``to be
submitted to the Dicastery for approval'' (art.~1); the possible
conclusions are six exactly---\term{Nihil obstat}, \term{Prae
oculis habeatur}, \term{Curatur}, \term{Sub mandato},
\term{Prohibetur et obstruatur}, \term{Declaratio de non
supernaturalitate} (nn.~17--22); and, as n.~23 states the
governing rule, ``as a rule, neither the Diocesan Bishop, nor
the Episcopal Conferences, nor the Dicastery will declare that
these phenomena are of supernatural origin, even if a \term{Nihil
obstat} is granted\ldots\ It remains true, however, that the
Holy Father can authorize a special procedure in this
regard.''\endnote{DDF, \work{Norms for Proceeding in the
Discernment of Alleged Supernatural Phenomena} (17 May 2024),
arts.~1--2, nn.~16--23, and art.~27 (replacement of the 1978
norms), English text (References), read 2026-07-25. A
\term{Nihil obstat} under n.~17 acknowledges ``signs of the
action of the Holy Spirit'' and permits and encourages pastoral
appreciation ``[w]ithout expressing any certainty about the
supernatural authenticity of the phenomenon itself.''}

The 2024 norms therefore mark a deliberate contrast with the
very genre of the 1862 act: the ordinary modern conclusion is a
pastoral \term{Nihil obstat} without a supernaturality
declaration, where Bishop Laurence, under the discipline of his
own time, judged and declared the apparition real and divine.
The norms do not erase that contrast retroactively.

\subsection{The 1862 judgment as a legacy act}

The 1862 mandement is a completed positive judgment of the
competent authority of its era, and it stands in its own
historical formula. Nothing in the 2024 norms relabels it: it is
not a \term{Nihil obstat}, and this study does not translate it
into one; it is what its own words say---\term{nous jugeons que
l'Immacul\'ee Marie\ldots\ a r\'eellement apparu}, with the
faithful \term{fond\'es \`a la croire certaine}---made under the
pre-1978 regime in which the diocesan bishop's judgment was the
ordinary conclusive act, and expressly submitted by its author
to the Roman Pontiff. The apparition it approved has since been
woven into the universal Church's liturgy, teaching, and
practice by the distinct acts of \S\S6--7, each with its own
object. No later act of the Holy See or of the diocese
superseding, restricting, or reversing the 1862 judgment was
located by this project; the modern Holy See documents that
touch Lourdes---from Pius X's 1907 feast extension to the 2024
norms' own presentation, which cites the Lourdes-anchored
explanation of approval---uniformly build on the judgment rather
than reopen it.\endnote{Bounded negative search, recorded with
its limits: the DDF's online documents index
(\url{https://www.doctrinafidei.va/en/documenti.html}, English
list, checked 2026-07-25) contains no Lourdes entry; web
searches in English and French for any DDF act under the 2024
norms concerning Lourdes returned none; the diocesan and
Sanctuary sites present the 1862 judgment as standing. The
search did not extend to print-only registers, and silence is
not a guarantee (Appendix~A).}

\subsection{Controlling-status dossier}

\begin{statusdossier}{Lourdes 1858 --- controlling status as of
2026-07-25}
\dosfield{Event identity}{Eighteen reported apparitions of the
Blessed Virgin Mary to Bernadette Soubirous, grotto of
Massabielle, Lourdes, 11 February--16 July 1858.}
\dosfield{Jurisdiction}{Diocese of Tarbes in 1858--1862; today
the Diocese of Tarbes and Lourdes (Bishop Jean-Marc Micas,
appointed 2022), Latin Church, France.}
\dosfield{Controlling act}{Mandement of Bishop Bertrand-S\'ev\`ere
Laurence, 18 January 1862: judgment that the Immaculate Mary,
Mother of God, really appeared eighteen times; the faithful
justified in believing it certain; diocesan cult authorized;
judgment submitted to the Roman Pontiff. Witness: Lasserre 1870
reproduction collated with Sanctuary transcriptions; originals
not consulted.}
\dosfield{Procedural regime}{Pre-1978 diocesan discernment,
following Benedict XIV's rules; episcopal commission of 28 July
1858; consulted physicians and scientists.}
\dosfield{Later acts (own objects)}{Seven cure recognitions in
the 1862 act and sixty-five further episcopal cure recognitions
through 16 April 2025 (total 72); proper Office and Mass 1890;
universal feast 13 November 1907 (11 February); Bureau des
Constatations 1883 with Roman assent 1886/1905; CMIL 1947/1954;
beatification 14 June 1925; canonization 8 December 1933; World
Day of the Sick instituted 13 May 1992; papal visits 1983, 2004,
2008 (announced: 27 September 2026).}
\dosfield{Current status}{Positive legacy diocesan judgment of
1862 in force in its historical formula; no superseding,
restricting, or reversing act located; no DDF act under the 2024
norms concerns Lourdes (DDF documents index checked 2026-07-25);
the 2024 norms govern any new discernment, not this completed
one.}
\dosfield{Not established by any act}{Obligation of divine
faith; verbal inerrancy of any reported sentence; scientific
demonstration of supernatural causation in any cure; curative
property of the water; authentication of every Lourdes
narrative, image, or publication; any doctrine beyond the
deposit of faith.}
\dosfield{Register}{``Reported'' remains the default register
for all words and phenomena attributed to the apparition; the
1862 act's own qualifications and this study's ceilings are in
\S4 and Appendix~A.}
\end{statusdossier}

The standing of Lourdes, stated in one sentence for the reader
who needs only one: a competent, completed, positively judged,
and continuously received private revelation, whose approval
authorizes prudent belief and warm devotion, obliges neither,
and leaves the whole weight of faith where it always was---on
the public Revelation of God in Jesus Christ, which Lourdes at
its best does nothing but serve.
